\documentclass{article}

\usepackage{amsmath}
\usepackage{amssymb}
\usepackage{xcolor}
\usepackage{tikz}
\usepackage[shortlabels]{enumitem}

\usepackage[utf8]{inputenc}
\usepackage[T1]{fontenc}

\usepackage[polish]{babel}

\usepackage{listings}
\lstset{
    language=SQL,
    inputencoding=utf8x,
    extendedchars=\true,
    literate={ą}{{\k{a}}}1
             {Ą}{{\k{A}}}1
             {ę}{{\k{e}}}1
             {Ę}{{\k{E}}}1
             {ó}{{\'o}}1
             {Ó}{{\'O}}1
             {ś}{{\'s}}1
             {Ś}{{\'S}}1
             {ł}{{\l{}}}1
             {Ł}{{\L{}}}1
             {ż}{{\.z}}1
             {Ż}{{\.Z}}1
             {ź}{{\'z}}1
             {Ź}{{\'Z}}1
             {ć}{{\'c}}1
             {Ć}{{\'C}}1
             {ń}{{\'n}}1
             {Ń}{{\'N}}1
}

\title{Dekompozycja macierzy}
\date{2026}

\begin{document}
\maketitle

\begin{enumerate}

\item Znaleźć rozkład SVD dla macierzy:
    \begin{enumerate}[a)]
        \item $A = \begin{bmatrix}\begin{array}{rrr} 1 & 2 & 0\\ 2 & 1 & 0\\ 0 & 0 & 1 \end{array}\end{bmatrix}$,
        \item $A = \begin{bmatrix}\begin{array}{rrr} 1 & 2 & 0\\ 2 & 0 & 2 \end{array}\end{bmatrix}$,
    \end{enumerate}

    Rozkład SVD (rozkład według wartości osobliwych - \textit{Singular Value Decomposition}) jset to rozkład macierzy $A \in \mathbb{R}^{m \times n}$ na trzy specyficzne macierze: $U \in \mathbb{R}^{m \times m}$ - ortogonalna, $\Sigma \in \mathbb{R}^{m \times n}$ - o postaci bliskiej do ortogonalnej, $V \in \mathbb{R}^{n \times n}$ - ortogonalna.
    $$
        A = U \Sigma V^{T}
    $$

\item Wyznaczyć rozkład Cholsekiego macierzy:
    \begin{enumerate}[a)]
        \item $A = \begin{bmatrix}\begin{array}{rr} 2 & 1 \\ 1 & 1 \end{array}\end{bmatrix}$,
        \item $A = \begin{bmatrix}\begin{array}{rrr} 2 & 1 & 0\\ 1 & 1 & 0\\ 0 & 0 & 1 \end{array}\end{bmatrix}$,
        \item $A = \begin{bmatrix}\begin{array}{rrr} 9 & 0 & 0\\ 0 & 4 & 2\\ 0 & 2 & 2 \end{array}\end{bmatrix}$.
    \end{enumerate}

    Rozkład Choleskiego jest procedurą rozkładu symetrycznej, dodatnio określonej macierzy $A \in \mathbb{R}^{n \times n}$ do iloczynu:
    $$
        A = LL^{T}
    $$
    gdzie macierz $L$ jest macierzą trójkątną dolną, zatem:
    $$
        \begin{bmatrix}\begin{array}{rrr} a_{11} & \cdots & a_{1n}\\ \vdots & \cdots & \vdots \\ a_{n1} & \cdots & a_{nn} \end{array}\end{bmatrix}
        =
        \begin{bmatrix}\begin{array}{rrrr} l_{11} & 0 & \cdots & 0\\ l_{21} & l_{22} & \vdots \\ \vdots & \cdots & \cdots & \vdots \\ l_{n1} & l_{n2} & \cdots & l_{nn} \end{array}\end{bmatrix}
        \begin{bmatrix}\begin{array}{rrrr} l_{11} & l_{21} & \cdots & l_{n1}\\ 0 & l_{22} & \vdots & l_{n2}\\ \vdots & \cdots & \cdots & \vdots \\ 0 & \cdots & 0 & l_{nn} \end{array}\end{bmatrix}
    $$
    Przy takiej notacji łatwo zuważyć, że:
    \begin{itemize}
        \item $a_{11} = l_{11}^{2} \rightarrow l_{11}=\sqrt{a_{11}}$,
        \item $a_{21} = l_{21}l_{11} \rightarrow l_{21} = \frac{a_{21}}{l_{11}}$,
        \item $\cdots$
    \end{itemize}
    A w ogólności:
    \begin{itemize}
        \item $l_{ii} = \sqrt{a_{ii} - \sum_{k=0}^{i-1}{l_{ik}^{2}}}$,
        \item $l_{ji} = \frac{a_{ji} - \sum_{k=0}^{i-1}{l_{jk}l_{ik}}}{l_{ii}}$
    \end{itemize}

\item Wyznaczyć rozkład $LU$ macierzy:
    \begin{enumerate}[a)]
        \item $A = \begin{bmatrix}\begin{array}{rr} 2 & 1 \\ 1 & 1 \end{array}\end{bmatrix}$,
        \item $A = \begin{bmatrix}\begin{array}{rrr} 1 & 2 & 0\\ 2 & 1 & 0\\ 0 & 0 & 1 \end{array}\end{bmatrix}$,
        \item $A = \begin{bmatrix}\begin{array}{rrr} 5 & 3 & 1\\ 1 & 2 & 0\\ 3 & 0 & 4 \end{array}\end{bmatrix}$,
        \item $A = \begin{bmatrix}\begin{array}{rrr} 2 & 1 & 1\\ 1 & 1 & 1\\ 0 & 0 & 4 \end{array}\end{bmatrix}$.
    \end{enumerate}

    Rozkład $LU$ jest procedurą rozkładu macierzy $A \in \mathbb{R}^{n \times n}$ do iloczynu:
    $$
        A = LU
    $$
    gdzie: $L \in \mathbb{R}^{n \times n}$ - macierz trójkątna dolna, $U \in \mathbb{R}^{n \times n}$ - macierz trójkątna górna.

    Dekompozycji $LU$ można dokonać metodą Doolittle'a - zakładając, że przekątna macierzy $L$ wypoełniona jest jedynkami,zatem:
    $$
        \begin{bmatrix}\begin{array}{rrr} a_{11} & \cdots & a_{1n}\\ \vdots & \cdots & \vdots \\ a_{n1} & \cdots & a_{nn} \end{array}\end{bmatrix}
        =
        \begin{bmatrix}\begin{array}{rrrr} 1 & 0 & \cdots & 0\\ l_{21} & 1 & \vdots \\ \vdots & \cdots & \cdots & \vdots \\ l_{n1} & l_{n2} & \cdots & 1 \end{array}\end{bmatrix}
        \begin{bmatrix}\begin{array}{rrrr} u_{11} & u_{12} & \cdots & u_{1n}\\ 0 & u_{22} & \vdots & u_{2n}\\ \vdots & \cdots & \cdots & \vdots \\ 0 & \cdots & 0 & u_{nn} \end{array}\end{bmatrix}
    $$
    Wówczas dla $i \in \{ 1, 2, \cdots, n \}$:
    \begin{itemize}
        \item $u_{ij} = a_{ij} - \sum_{k=1}^{i-1}{l_{ik}u_{kj}}, j\in {i, i+1, \cdots, n}$,
        \item $l_{ji} = \frac{1}{u_{ii}}(a_{ji} - \sum_{k=1}^{i-1}{l_{jk}u_{ki}}), j\in {i+1, i+2, \cdots, n}$
    \end{itemize}

\end{enumerate}
    

\newpage
\textbf{Odpowiedzi:}
\begin{description}
    \item [1:] a) $U = \begin{bmatrix}\begin{array}{ccc} -\frac{\sqrt(2)}{2} & 0 & -\frac{\sqrt(2)}{2}\\ -\frac{\sqrt(2)}{2} & 0 & \frac{\sqrt(2)}{2}\\ 0 & 1 & 0 \end{array}\end{bmatrix}$, $\Sigma = \begin{bmatrix}\begin{array}{ccc} 3 & 0 & 0\\ 0 & 1 & 0\\ 0 & 0 & 1 \end{array}\end{bmatrix}$, $V = \begin{bmatrix}\begin{array}{ccc} -\frac{\sqrt(2)}{2} & 0 & \frac{\sqrt(2)}{2}\\ -\frac{\sqrt(2)}{2} & 0 & -\frac{\sqrt(2)}{2}\\ 0 & 1 & 0 \end{array}\end{bmatrix}$,
    b) $U = \begin{bmatrix}\begin{array}{cc} \frac{1}{\sqrt(5)} & \frac{2}{\sqrt(5)}\\ \frac{2}{\sqrt(5)} & -\frac{1}{\sqrt(5)} \end{array}\end{bmatrix}$, $\Sigma = \begin{bmatrix}\begin{array}{ccc} 3 & 0 & 0\\ 0 & 2 & 0 \end{array}\end{bmatrix}$, $V = \begin{bmatrix}\begin{array}{rrr} \frac{5}{3\sqrt(5)} & 0 & -\frac{2}{3}\\ \frac{2}{3\sqrt(5)} & \frac{2}{\sqrt(5)} & \frac{1}{3}\\ \frac{4}{3\sqrt(5)} & -\frac{1}{\sqrt(5)} & \frac{2}{3} \end{array}\end{bmatrix}$

    \item [2:] a)  $L = \begin{bmatrix}\begin{array}{cc} \sqrt(2) & \frac{\sqrt(2)}{2}\\ 0 & \frac{\sqrt(2)}{2}\\ \end{array}\end{bmatrix}$
    b) $L = \begin{bmatrix}\begin{array}{ccc} \sqrt(2) & \frac{\sqrt(2)}{2} & 0\\ 0 & \frac{\sqrt(2)}{2} & 0\\ 0 & 0 & 1 \end{array}\end{bmatrix}$,
    b) $L = \begin{bmatrix}\begin{array}{ccc} 3 & 0 & 0\\ 0 & 2 & 1\\ 0 & 0 & 1 \end{array}\end{bmatrix}$
     

    \item [3:] a) $L = \begin{bmatrix}\begin{array}{rr} 1 & 0\\ \frac{1}{2} & 1\\ \end{array}\end{bmatrix}$, $U = \begin{bmatrix}\begin{array}{rr} 2 & 1\\ 0 & \frac{1}{2} \end{array}\end{bmatrix}$,
    b) $L = \begin{bmatrix}\begin{array}{rrr} 1 & 0 & 0\\ \frac{1}{2} & 1 & 0\\ 0 & 0 & 1 \end{array}\end{bmatrix}$, $U = \begin{bmatrix}\begin{array}{rrr} 2 & 1 & 0\\ 0 & \frac{3}{2} & 0\\ 0 & 0 & 1 \end{array}\end{bmatrix}$,
    c) $L = \begin{bmatrix}\begin{array}{rrr} 5 & 0 & 0\\ \frac{1}{5} & 1 & 0\\ \frac{3}{5} & -\frac{9}{7} & 1 \end{array}\end{bmatrix}$, $U = \begin{bmatrix}\begin{array}{rrr} 5 & 3 & 2\\ 0 & \frac{7}{5} & -\frac{2}{5}\\ 0 & 0 & \frac{16}{7} \end{array}\end{bmatrix}$,
    d) $L = \begin{bmatrix}\begin{array}{rrr} 1 & 0 & 0\\ \frac{1}{2} & 1 & 0\\ 0 & 0 & 1 \end{array}\end{bmatrix}$, $U = \begin{bmatrix}\begin{array}{rrr} 2 & 1 & 1\\ 0 & \frac{1}{2} & \frac{1}{2}\\ 0 & 0 & 4 \end{array}\end{bmatrix}$,
\end{description}

\end{document}